\documentclass[12pt]{article}

\usepackage[margin = .8in]{geometry}
\usepackage{amsmath}
\usepackage{graphicx}
\usepackage{multicol, enumerate, tabularx}

\usepackage[parfill]{parskip}

\usepackage{adjustbox, soul}

\usepackage{fancyhdr}
\pagestyle{fancy}

\lhead{\emph{Math F113X: Loans with regular payments }}
\rhead{\emph{lecture notes}}

\renewcommand{\emph}[1]{\textsf{\textbf{#1}}}
\usepackage{tikz}
\usetikzlibrary{calc,trees,positioning,arrows,fit,shapes,through, backgrounds}
\usetikzlibrary{patterns}

\usetikzlibrary{decorations.markings}
\usetikzlibrary{arrows}

\usepackage{pgfplots}

\usepackage{longtable}
\usepackage{tabularx}

\newcommand{\ds}{\displaystyle}
\newcommand{\ans}[1][1in]{\rule{#1}{.5pt}}

\newcommand{\points}[1]{(#1 points.)}		% Trying to be lazy.

\usepackage{array}
\newcolumntype{L}[1]{>{\raggedright\let\newline\\\arraybackslash\hspace{0pt}}m{#1}}
\newcolumntype{C}[1]{>{\centering\let\newline\\\arraybackslash\hspace{0pt}}m{#1}}
\newcolumntype{R}[1]{>{\raggedleft\let\newline\\\arraybackslash\hspace{0pt}}m{#1}}
\newcommand{\red}[1]{\textcolor{red}{#1}}

\newcommand{\be}{\begin{enumerate}}
\newcommand{\ee}{\end{enumerate}}

\begin{document}
%\begin{center}
%{\large  Finance Section 4: Introduction to Loan Calculation}
%\end{center}
Goal: Think about loan calculations intuitively. What happens if you make regular payments?

\begin{enumerate}
\item  (Basic Example:) Suppose you take out a loan for \$10,000 with an annual interest rate of 5\% compounded annually, and you will make payments annually.
	\begin{enumerate}
	\item What is the \emph{minimum} annual payment that ensures you will eventually pay off the loan? (Round your answer to the nearest dollar.)
	\vfill
	\item Suppose you decided to make annual payments of \$600. How much do you owe at the end of one year? (One year of interest is added and one year of payments is subtracted.)
	\vfill
	\item How much do you owe at the end of the second year? The third year?
	\vfill
	\item Use a \emph{spreadsheet} to determine how long it will take to pay off the loan. (Be thoughtful about how you set this up so you can change the parameters!)
	\vfill
	\item How much did the loan cost you? How much was interest?
	\vfill
	\item What happens if you increase your annual payments to \$800?
	\vfill
	\end{enumerate}
\newpage
\item (Compounded Monthly Example:) Suppose you take out a loan for \$10,000 with an annual interest rate of 5\% compounded \emph{monthly} and you are going to make payments monthly.
	\begin{enumerate}
	\item What is the \emph{minimum} monthly payment that ensures you will eventually pay off the loan? (Round your answer to the nearest dollar.)
	\vfill
	\item If you make a \$200 payment each month, how many \emph{years} will it take to pay off the loan and how much did you pay in total?
	\vfill
	\item If you make a \$299.71 payment each month, how many \emph{years} will it take to pay off the loan?
	\vfill
	
	\item \emph{Formulas:}
	
	\begin{tabular}{| l | l |}
\hline
$P$ = principal / starting amount & $r$ = annual interest rate (APR) \\ \hline
$I$ = interest & $n$ = number of compounding periods per year\\ \hline
$A$ = final amount & $t$ = length of the loan, in years\\ \hline
& $d$ = regular loan payment \\ \hline
\end{tabular}


\begin{center}
\begin{tabularx}{.8\linewidth}{|X | X|} 
\hline
Loan amount given payment & payment given loan amount\\ \hline
$\displaystyle P = \frac{d\left(1 - \left(1+\frac{r}{n}\right)^{(-nt)}\right)}{\left(\frac{r}{n}\right)}$ & $\displaystyle d = \frac{P\left(\frac{r}{n}\right)}{\left(1 - \left(1+\frac{r}{n}\right)^{(-nt)}\right)}$\\
\hline
\end{tabularx}
\end{center}


As before, suppose you take out a loan for \$10,000 with an annual interest rate of 5\% compounded \emph{monthly} and you are going to make payments monthly. If you want to pay off the loan in 3 years, what should your monthly payment be?

\vfill




	
	
	\end{enumerate}


\end{enumerate}
\end{document}
